\section*{Chapitre \ref{ch:g7:volcanoes} --- Les volcans}

\begin{solution}{exo:g7:volcanoes:1}
(a) \omterm{def:g7:volcanoes:vent}{cratère}~; (b) \omterm{def:g7:volcanoes:magma}{magma}~; (c) \omterm{def:g7:volcanoes:vent}{bouche éruptive}~; (d) \omterm{def:g4:volcanoes-earthquakes-first-look:eruption}{lave}~; (e) \omterm{def:g7:volcanoes:magma-chamber}{chambre magmatique}.
\end{solution}

\begin{solution}{exo:g7:volcanoes:2}
(a) Effusive. (b) Explosive. (c) Effusive~: une fontaine de \omterm{def:g4:volcanoes-earthquakes-first-look:eruption}{lave} projette
des gouttes incandescentes qui retombent et s'écoulent~; le \omterm{def:g7:volcanoes:magma}{magma} n'est pas
réduit en cendres. (d) Explosive~: c'est une \omterm{def:g7:volcanoes:pyroclastic-flow}{coulée pyroclastique}.
\end{solution}

\begin{solution}{exo:g7:volcanoes:3}
(a) Vrai. (b) Faux~: le gaz volcanique le plus abondant est la vapeur
d'eau~; le dioxyde de carbone vient ensuite. (c) Faux~: une \omterm{def:g7:volcanoes:pyroclastic-flow}{coulée
pyroclastique} est une avalanche de gaz brûlants, de cendres et de morceaux
de \omterm{def:g3:stones-pebbles-sand:rock}{roche}. (d) Faux~: ils se trouvent à l'intérieur de la \omterm{def:g7:plate-tectonics:plate}{plaque} pacifique,
à plus de 3200~km de la limite la plus proche, au-dessus d'un \omterm{def:g7:volcanoes:hotspot}{point chaud}.
\end{solution}

\begin{solution}{exo:g7:volcanoes:4}
(a) Environ 1000 à 1200~degrés. (b) Environ 600 à 800~degrés. (c) Près de
la \omterm{def:g7:volcanoes:vent}{bouche}, où la \omterm{def:g4:volcanoes-earthquakes-first-look:eruption}{lave} est jaune (1000 à 1200~degrés) plutôt qu'orange (800
à 1000~degrés)~: la \omterm{def:g4:volcanoes-earthquakes-first-look:eruption}{lave} se refroidit en s'éloignant de la \omterm{def:g7:volcanoes:vent}{bouche}.
\end{solution}

\begin{solution}{exo:g7:volcanoes:5}
(a) $838 + 221 = 1059$ \omterm{def:g4:volcanoes-earthquakes-first-look:volcano}{volcans}. (b) $1059 \div 1214 \approx 0.872$~:
environ \textbf{87~\%}. (c) $69 + 18 + 12 + 0.5 = 99.5$, environ
100~\%~: le petit écart vient de l'arrondi de chaque pourcentage.
\end{solution}

\begin{solution}{exo:g7:volcanoes:6}
(a) $564 \div 660 \approx 0.855$~: environ \textbf{85~\%}. (b)
$660 - 564 - 62 - 33 = 1$~: un seul cône raide. (c) Les cônes raides sont
construits par une \omterm{def:g4:volcanoes-earthquakes-first-look:eruption}{lave} visqueuse et des cendres, la marque des \omterm{def:g7:volcanoes:eruption-types}{éruptions
explosives}~: le long de la ceinture de feu, là où une \omterm{def:g7:plate-tectonics:plate}{plaque} s'enfonce
sous une autre, on s'attend à de nombreuses \omterm{def:g7:volcanoes:eruption-types}{éruptions explosives}, comme
celles du mont Saint Helens et du Pinatubo.
\end{solution}

\begin{solution}{exo:g7:volcanoes:7}
(a) $1170 - 1140 = 30$~degrés. (b) $30 \div 10 = 3$~degrés par
kilomètre. (c) Dans le tunnel, la \omterm{def:g4:volcanoes-earthquakes-first-look:eruption}{lave} ne touche pas l'air~: le toit durci
qui la recouvre garde la chaleur, comme le couvercle d'une casserole. À
l'air libre, sa surface cède aussitôt sa chaleur à l'air.
\end{solution}

\begin{solution}{exo:g7:volcanoes:8}
(a) Kauai, d'au moins $5.5 - 0.7 = 4.8$~millions d'années. (b) Kauai~: la
\omterm{def:g4:clouds-rain-snow:precipitation}{pluie} et les vagues l'usent depuis bien plus longtemps. (c) Big Island, où
de la \omterm{def:g4:volcanoes-earthquakes-first-look:eruption}{lave} nouvelle sort encore. (d) Vers le nord-ouest~: Kauai, emportée
la première loin du \omterm{def:g7:volcanoes:hotspot}{point chaud}, se trouve maintenant au nord-ouest de
celui-ci.
\end{solution}

\begin{solution}{exo:g7:volcanoes:9}
(a) Presque tous les \omterm{def:g4:volcanoes-earthquakes-first-look:volcano}{volcans} se trouvent sur des \omterm{def:g7:plate-tectonics:plate-boundary}{limites de plaques}
(87~\% dans la liste)~: des \omterm{def:g4:volcanoes-earthquakes-first-look:volcano}{volcans} aussi loin de toute limite sont une
exception. (b) Profond dans le \omterm{def:g6:inside-the-earth:layers}{manteau}, sous Hawaï, un \omterm{def:g7:volcanoes:hotspot}{point chaud} reste
chaud pendant des millions d'années. Il fait fondre en partie la \omterm{def:g3:stones-pebbles-sand:rock}{roche}
située au-dessus. Le \omterm{def:g7:volcanoes:magma}{magma}, plus léger que la \omterm{def:g3:stones-pebbles-sand:rock}{roche} solide qui l'entoure,
monte à travers la \omterm{def:g7:plate-tectonics:plate}{plaque} et sort en \omterm{def:g4:volcanoes-earthquakes-first-look:eruption}{éruption}. Pendant que la \omterm{def:g7:plate-tectonics:plate}{plaque} se
déplace, de nouveaux \omterm{def:g4:volcanoes-earthquakes-first-look:volcano}{volcans} grandissent au-dessus du \omterm{def:g7:volcanoes:hotspot}{point chaud} et les
anciens sont emportés. (c) Des points orange («~à l'intérieur d'une
\omterm{def:g7:plate-tectonics:plate}{plaque}~»).
\end{solution}

\begin{solution}{exo:g7:volcanoes:10}
Le premier \omterm{def:g7:volcanoes:magma}{magma} avait perdu la plus grande partie de ses gaz en montant,
comme une bouteille d'eau gazeuse restée ouverte ou restée posée~: en
arrivant à la surface, il ne lui restait rien pour le faire éclater, et il
est sorti lentement. Le \omterm{def:g7:volcanoes:magma}{magma} des 12 et 15~juin contenait encore ses gaz,
comme la bouteille secouée~: quand la pression a baissé près de la
surface, les gaz ont formé des bulles d'un seul coup, qui ont gonflé et
déchiqueté le \omterm{def:g7:volcanoes:magma}{magma} visqueux en cendres. Le premier \omterm{def:g7:volcanoes:magma}{magma} est celui qui
avait perdu ses gaz.
\end{solution}

\begin{solution}{exo:g7:volcanoes:11}
(a) $10 \div 80 = 0.125$~heure, et $0.125 \times 60 = 7.5$~minutes. (b)
$10 \div 130 \approx 0.077$~heure, soit environ $0.077 \times 60 \approx
4.6$~minutes, donc environ 5~minutes. (c) La famille a besoin de
20~minutes, mais la coulée arrive en 5 à 8~minutes~: quand on voit la
coulée, il est trop tard. La famille n'est en sécurité que si elle part
quand les scientifiques voient les signes avant-coureurs, avant
l'\omterm{def:g4:volcanoes-earthquakes-first-look:eruption}{éruption}.
\end{solution}

\begin{solution}{exo:g7:volcanoes:12}
(a) $111 \times 360 = 39\,960$~: environ \textbf{40\,000~km}. (b)
$40\,000 \div 15 \approx 2667$~: environ 2\,670~km par jour. (c)
$2667 \div 24 \approx 111$~: environ \textbf{110~km par heure}. (d) Par le
vent~: les cendres fines n'ont plus de vitesse propre une fois que
l'\omterm{def:g4:volcanoes-earthquakes-first-look:eruption}{éruption} les a projetées en l'air~; elles flottent haut dans l'air et le
vent les emporte, tandis qu'elles retombent lentement.
\end{solution}

\begin{solution}{pb:g7:volcanoes:1}
\textbf{1.} Du 16 au 31~mars~: $31 - 16 = 15$~jours~; puis 30~jours en
avril et 18~jours en mai~: $15 + 30 + 18 = 63$~jours.

\textbf{2.} $10\,000 \div 62 \approx 161$~: plus de 160 petits \omterm{def:g4:volcanoes-earthquakes-first-look:earthquake}{séismes} par
jour, en moyenne.

\textbf{3.} Le gonflement montrait que du \omterm{def:g7:volcanoes:magma}{magma} était monté haut dans la
\omterm{def:g1:landscapes-around-us:landform}{montagne} et poussait la \omterm{def:g3:stones-pebbles-sand:rock}{roche} vers l'extérieur. Il rendait aussi le flanc
nord raide et instable, prêt à glisser.

\textbf{4.} Un triangle bleu se trouve là où une \omterm{def:g7:plate-tectonics:plate}{plaque} s'enfonce sous une
autre. On s'attend à un \omterm{def:g7:volcanoes:magma}{magma} visqueux et riche en gaz, à des \omterm{def:g7:volcanoes:eruption-types}{éruptions
explosives} et à un cône raide (un stratovolcan) --- exactement ce que
montre la première photographie.

\textbf{5.} $240 \div 60 = 4$~kilomètres par minute.

\textbf{6.} Le flanc gonflé retenait le \omterm{def:g7:volcanoes:magma}{magma} sous pression, comme le
bouchon d'une bouteille. Quand il a glissé, la pression a chuté d'un
coup~; les gaz dissous dans le \omterm{def:g7:volcanoes:magma}{magma} se sont transformés en bulles tous à
la fois et ont jailli par le flanc ouvert --- de côté, là où le couvercle
avait disparu.

\textbf{7.} $27 \div 480 \approx 0.056$~heure, et $0.056 \times 60 \approx
3.4$~minutes. Comme le souffle avançait au moins aussi vite, il a mis au
plus environ 3.4~minutes.

\textbf{8.} $8 \div 130 \approx 0.062$~heure, et $0.062 \times 60 \approx
3.7$~: environ \textbf{4~minutes}.

\textbf{9.} $24 \div 8.8 \approx 2.7$~: la colonne est montée environ
2.7~fois plus haut que le mont Everest.

\textbf{10.} Effusives~: la \omterm{def:g4:volcanoes-earthquakes-first-look:eruption}{lave} visqueuse est sortie sans exploser et
s'est entassée au-dessus de la \omterm{def:g7:volcanoes:vent}{bouche}. Comme la bouteille restée posée, ce
\omterm{def:g7:volcanoes:magma}{magma} avait déjà perdu la plus grande partie de ses gaz.

\textbf{11.} Avant, la \omterm{def:g1:landscapes-around-us:landform}{montagne} était un cône lisse et pointu couvert de
\omterm{def:g4:clouds-rain-snow:precipitation}{neige}, au-dessus de forêts vertes. Après, son sommet et son flanc nord ont
disparu~: un immense \omterm{def:g7:volcanoes:vent}{cratère} ouvert vers nous les a remplacés, et les
pentes et la \omterm{def:g1:landscapes-around-us:landform}{vallée} au premier plan sont grises de cendres et de \omterm{def:g3:stones-pebbles-sand:rock}{roches}
brisées.

\textbf{12.} $2950 - 2549 = 401$~: la \omterm{def:g1:landscapes-around-us:landform}{montagne} a perdu environ
\textbf{401~m} de hauteur, soit environ 400~mètres, comme on l'a vu en
année~4.
\end{solution}
